\chapter{Introduction}

\section{Background and Motivation}

The rapid evolution of Large Language Models (LLMs) such as GPT-4, Gemini, and LLaMA has revolutionized natural language processing, enabling systems that can generate human-like text, summarize documents, and answer complex questions. However, a fundamental limitation of standalone LLMs is their tendency to \textit{hallucinate}---generating plausible-sounding but factually incorrect information, particularly when queried about domain-specific or time-sensitive topics~\cite{lewis2020retrieval}.

Retrieval-Augmented Generation (RAG) addresses this limitation by augmenting the LLM's generative capabilities with an external knowledge retrieval step. In a RAG pipeline, user queries are first converted into vector embeddings, which are then used to search a vector database for semantically relevant document chunks. These retrieved chunks serve as context for the LLM, grounding its response in verified factual data.

While RAG significantly improves response accuracy and traceability, it introduces new attack surfaces that are particularly dangerous in sensitive domains such as banking and financial services:

\begin{enumerate}
    \item \textbf{Prompt Injection Attacks}: Adversaries craft malicious inputs that manipulate the LLM into ignoring its system instructions, revealing confidential data, or executing unintended operations~\cite{perez2022prompt}.
    
    \item \textbf{Embedding Inversion Attacks}: Attackers with access to stored vector embeddings can reverse-engineer them to reconstruct the original text, potentially exposing sensitive PII such as account numbers, phone numbers, and financial transaction details~\cite{song2020information}.
\end{enumerate}

\section{Problem Statement}

Given a banking enterprise deploying a RAG-based AI assistant for customer queries, the challenge is to design a system that:

\begin{itemize}
    \item Processes PDF documents containing sensitive financial data through a secure pipeline before vectorization.
    \item Detects and neutralizes prompt injection, SQL injection, cross-site scripting (XSS), and shell command injection attacks in real-time.
    \item Prevents embedding inversion attacks through mathematically provable differential privacy mechanisms.
    \item Ensures multi-tenant data isolation so that one customer's documents are invisible to other customers.
    \item Provides controlled de-pseudonymization where only the authenticated document owner can recover original PII values.
\end{itemize}

\section{Objectives}

The primary objectives of this project are:

\begin{enumerate}
    \item To design and implement a \textbf{7-layer security pipeline} for document ingestion in a RAG system.
    
    \item To develop a \textbf{context-aware PII detection engine} using a 25-character look-behind algorithm specific to Indian financial document formats.
    
    \item To implement \textbf{AES-256-CBC field-level encryption} for pseudonym mappings stored in MongoDB.
    
    \item To apply \textbf{$\epsilon$-differential privacy} through Laplacian noise injection on 3072-dimensional Gemini embeddings to prevent embedding inversion.
\end{enumerate}

\section{Scope of the Project}

The system is designed as a RAG application targeting internal bank operations. The scope includes:

\begin{itemize}
    \item A backend service layer deployed on cloud infrastructure with secured NoSQL and vector databases.
    \item A client application interface with role-specific dashboards for administrative and customer operations.
    \item Processing of PDF documents containing Indian banking formats (IFSC codes, 9--18 digit account numbers, UPI IDs).
    \item Real-time threat detection, PII masking, and differential privacy during document upload and query operations.
\end{itemize}

\section{Organization of the Report}

The remainder of this report is organized as follows:

\begin{itemize}
    \item \textbf{Chapter 2}: Literature Review --- surveys related work in RAG security, prompt injection defenses, embedding privacy, and PII protection techniques.
    \item \textbf{Chapter 3}: System Architecture --- presents the overall architecture, the 9-step document upload pipeline, and the 8-step query pipeline with mathematical foundations.
    \item \textbf{Chapter 4}: Implementation --- details the algorithms, code structure, encryption mechanisms, and differential privacy mathematics.
    \item \textbf{Chapter 5}: Results and Analysis --- demonstrates attack scenarios, defense outcomes, and system performance metrics.
    \item \textbf{Chapter 6}: Conclusion and Future Work --- summarizes contributions and outlines future enhancements.
\end{itemize}
